% Generated draft: requires mathematical and source review.
\documentclass[11pt]{article}
\usepackage[margin=25mm]{geometry}
\usepackage{fontspec}
\setmainfont{Latin Modern Roman}
\usepackage{amsmath,amssymb,mathtools,unicode-math}
\setmathfont{Latin Modern Math}
\usepackage{enumitem,needspace,xurl,hyperref}
\setlength{\parindent}{0pt}
\setlength{\parskip}{5pt}
\setlength{\emergencystretch}{3em}
\newcommand{\N}{\mathbb{N}}
\newcommand{\Z}{\mathbb{Z}}
\newcommand{\Q}{\mathbb{Q}}
\newcommand{\R}{\mathbb{R}}
\newcommand{\C}{\mathbb{C}}
\newcommand{\F}{\mathbb{F}}
\newcommand{\A}{\mathbb{A}}
\newcommand{\PP}{\mathbb{P}}
\newcommand{\E}{\mathbb{E}}
\newcommand{\eps}{\varepsilon}
\newcommand{\dd}{\,\mathrm{d}}
\begin{document}
\section*{UnsolvedMath ID: 701407}
\textbf{Source problem code:} catalogue entry
\subsection{Definitions and mathematical statement}
Let \(p\) be a prime. A pro-\(p\) group is a topological group isomorphic to an inverse limit of finite \(p\)-groups. Free pro-\(p\) groups and free pro-\(p\) products are understood in the category of pro-\(p\) groups with continuous homomorphisms.

For a pro-\(p\) group \(G\) and a positive integer \(n\), let \(F_p(x_1,\ldots,x_n)\) be the free pro-\(p\) group on the indicated finite basis, and put
\[
G\amalg_p F_p(x_1,\ldots,x_n)
\]
for its free pro-\(p\) product with \(G\). An equation over \(G\) is an expression \(v=1\) with \(v\) in this free pro-\(p\) product. For \(a=(a_1,\ldots,a_n)\in G^n\), let \(v(a)\) be the image of \(v\) under the unique continuous homomorphism to \(G\) that restricts to the identity on \(G\) and sends each \(x_i\) to \(a_i\). For any set \(S\) of such elements, define
\[
V_G(S)=\{a\in G^n: v(a)=1\text{ for every }v\in S\}.
\]
Thus solutions are sought in the Cartesian power \(G^n\), called the affine space in the source.

The group \(G\) is \emph{equationally Noetherian}, in the sense specified by this problem, if for every positive integer \(n\) and every set
\(S\subseteq G\amalg_p F_p(x_1,\ldots,x_n)\), there exists a finite subset \(S_0\subseteq S\) such that
\[
V_G(S)=V_G(S_0).
\]
Here \(V_G(\varnothing)=G^n\), and equivalence of systems means equality of their solution sets in \(G^n\).

\textbf{Problem (Kourovka Notebook 17.105).} Is every free pro-\(p\) group equationally Noetherian in this sense? The prime \(p\) is arbitrary. The supplied statement gives no restriction on the rank of the free pro-\(p\) group; whether a finite-rank restriction was intended remains unresolved. The equations above are elements of the free pro-\(p\) product, as explicitly required by the source.
\subsection{Short English statement}
Must every system of equations in finitely many variables over a free pro-\(p\) group have exactly the same solutions as some finite subsystem? Equations are interpreted in the free pro-\(p\) product, and solutions lie in the given group. The supplied source does not specify a rank restriction.
\subsection{Sources}
\begin{itemize}
\item Local UnsolvedMath registry, permanent ID \(701407\), exact source problem code \texttt{catalogue entry}, titled \emph{Kourovka Notebook Problem 17.105}. The supplied record labels the problem open; this label has not been independently verified.
\item \emph{The Kourovka Notebook (Unsolved Problems in Group Theory)}, original problem 17.105, 17th issue (2010); the supplied background also identifies the 21st issue edition (2026), \texttt{arXiv:1401.0300v48}. These bibliographic details are supplied source attribution and have not been independently checked. No source URL or page number is supplied.
\item Proposed in 2010 by N. S. Romanovski, with a truncated TeX diacritic at the end of the supplied proposer field. The exact spelling requires review. The collection description names E. I. Khukhro and V. D. Mazurov as editors.
\end{itemize}
\end{document}
